% The pipeline plate: one A1 sheet showing a map being made, stage by stage.
% Built with pdflatex. Shares the paper's palette and fonts.
\documentclass[a1paper,landscape]{article}

% A1 is four times A3 by area, so the body type is scaled to match: a poster set
% at 10pt on this sheet reads as a footnote from two metres away.
\usepackage[fontsize=18pt]{fontsize}

\usepackage[T1]{fontenc}
\usepackage[utf8]{inputenc}
\usepackage{libertinus-type1}
\usepackage[scaled=0.82]{beramono}
\usepackage[a1paper,landscape,top=40mm,bottom=34mm,left=40mm,right=40mm]{geometry}
\usepackage{graphicx}
\usepackage{xcolor}
\usepackage{booktabs}
\usepackage{tabularx}
\usepackage{array}
\usepackage{tikz}
\usepackage{enumitem}
\usetikzlibrary{positioning,arrows.meta,calc,fit}

\definecolor{ink}       {HTML}{1A1D21}
\definecolor{ink60}     {HTML}{5B6169}
\definecolor{ink30}     {HTML}{9AA1A9}
\definecolor{rule}      {HTML}{D8DBDF}
\definecolor{wash}      {HTML}{F4F5F7}
\definecolor{accent}    {HTML}{1F5C8B}
\definecolor{accentsoft}{HTML}{DCE8F1}
\definecolor{studioC}   {HTML}{0072B2}
\definecolor{agentC}   {HTML}{D55E00}

\newcolumntype{Y}{>{\raggedright\arraybackslash}X}
\newcolumntype{L}[1]{>{\raggedright\arraybackslash}p{#1}}
\newcommand{\id}[1]{{\ttfamily #1}}
\newcommand{\studio}[1]{\textcolor{studioC}{\textbf{#1}}}
\newcommand{\agent}[1]{\textcolor{agentC}{\textbf{#1}}}

\pagestyle{empty}
\setlength{\parindent}{0pt}
\renewcommand{\arraystretch}{1.22}

% A stage card: the render, its number, its name, and one line of what it added.
\newcommand{\stagecard}[4]{%
  \begin{minipage}[t]{0.132\textwidth}
    \fbox{\includegraphics[width=\linewidth]{#1}}\\[5pt]
    {\sffamily\small\color{ink30} #2}\ \ {\sffamily\small\bfseries\color{ink} #3}\\[2pt]
    {\footnotesize\color{ink60} #4}
  \end{minipage}%
}

% A stage card with no picture yet, so the sheet lays out before the renders land.
\newcommand{\stageplaceholder}[3]{%
  \begin{minipage}[t]{0.132\textwidth}
    \fcolorbox{rule}{wash}{\parbox[c][88mm][c]{\dimexpr\linewidth-2\fboxsep-2\fboxrule\relax}{\centering\color{ink30}\footnotesize stage render}}\\[5pt]
    {\sffamily\small\color{ink30} #1}\ \ {\sffamily\small\bfseries\color{ink} #2}\\[2pt]
    {\footnotesize\color{ink60} #3}
  \end{minipage}%
}

\begin{document}

% ---------------------------------------------------------------- masthead ----
{\sffamily\fontsize{92}{98}\selectfont\bfseries\color{ink} How a PGM Map Gets Made}\\[14pt]
{\sffamily\fontsize{34}{40}\selectfont\color{ink60}
 A map is made by two parties, and the line between them moves at every stage}\\[14pt]
{\color{rule}\rule{\textwidth}{1.6pt}}\\[12pt]

% ------------------------------------------------------------ stage ladder ----
{\sffamily\Large\bfseries\color{ink} One board, seven stages}\\[4pt]
{\large\color{ink60} Each picture is the same map exported again with one more
instrument switched on. Nothing between them is drawn by hand.}\\[12pt]

\IfFileExists{../figures/generated/stages/01-ground.png}{%
  \stagecard{../figures/generated/stages/01-ground.png}{1}{Ground}{The plan compiled. Flat, grey, rectangular.}\hfill
  \stagecard{../figures/generated/stages/02-relief.png}{2}{Relief}{Five marks and three pushes carve the fell.}\hfill
  \stagecard{../figures/generated/stages/03-paint.png}{3}{Paint}{A theme per bucket, banded by the ground's angle.}\hfill
  \stagecard{../figures/generated/stages/04-shapes.png}{4}{Shapes}{Fifteen authored, and the coast edited and bent.}\hfill
  \stagecard{../figures/generated/stages/05-layers.png}{5}{Layers}{Two made layers: a kiln's base and its stack.}\hfill
  \stagecard{../figures/generated/stages/06-rooms.png}{6}{Rooms}{The shells a spawn and a wool are stamped as.}\hfill
  \stagecard{../figures/generated/stages/07-dressed.png}{7}{Dressed}{Ten styles and nineteen props, each one seated.}%
}{%
  \stageplaceholder{1}{Ground}{The plan compiled. Flat, grey, rectangular.}\hfill
  \stageplaceholder{2}{Relief}{Five marks and three pushes carve the fell.}\hfill
  \stageplaceholder{3}{Paint}{A theme per bucket, banded by the ground's angle.}\hfill
  \stageplaceholder{4}{Shapes}{Fifteen authored, and the coast edited and bent.}\hfill
  \stageplaceholder{5}{Layers}{Two made layers: a kiln's base and its stack.}\hfill
  \stageplaceholder{6}{Rooms}{The shells a spawn and a wool are stamped as.}\hfill
  \stageplaceholder{7}{Dressed}{Ten styles and nineteen props, each one seated.}%
}

\vspace{10pt}
{\color{rule}\rule{\textwidth}{1pt}}\par\vspace{10pt}

% ----------------------------------------------------------- lower register ---
\noindent\begin{minipage}[t]{\textwidth}
\noindent
\begin{minipage}[t]{0.305\textwidth}
{\sffamily\Large\bfseries\color{ink} A map is four documents}\\[8pt]
{\color{ink60} Not four drafts. Four descriptions of the whole map at four grains,
each owned by a different tool.}\\[9pt]

\begin{tabularx}{\linewidth}{L{74mm} Y}
\toprule
\textbf{Level} & \textbf{What it states} \\
\midrule
The board & rectangles on a coarse grid, what each is for, where the objectives sit \\
The ground & real geometry at block resolution: outlines, heights, relief, paint \\
The play & teams, spawns, build regions, objectives and how they are captured \\
The map & the voxel world and the \id{map.xml} a server loads \\
\bottomrule
\end{tabularx}

\vspace{14pt}
{\color{ink60} The flow is one way. Nothing reads back up, and a finished map cannot
be decompiled into the plan that would have produced it.}
\end{minipage}
\hfill
\begin{minipage}[t]{0.305\textwidth}
{\sffamily\Large\bfseries\color{ink} Who decided this?}\\[8pt]
{\color{ink60} The question the whole paper asks, stage by stage.
\studio{Blue} is the studio. \agent{Orange} is the agent driving it.}\\[9pt]

\begin{tabularx}{\linewidth}{L{62mm} Y Y}
\toprule
\textbf{Stage} & \textbf{\studio{Studio}} & \textbf{\agent{Agent}} \\
\midrule
Arrangement & hub, spawn, wool approaches, frontline, land budget & which board, and everything the composer does not state \\
Elevation & the solved field & the marks, the grain, the pushes \\
Paint & which block on which cell & the theme, the axis, where the bands cut \\
Structures & storey heights, roof pitch, door side & which style, which seat from the offered set \\
Dressing & the legal seats & which seats to take \\
Objectives & regions, filters, monuments & what the map is played for \\
Refusal & what may not be built & what to do about it \\
\bottomrule
\end{tabularx}
\end{minipage}
\hfill
\begin{minipage}[t]{0.305\textwidth}
{\sffamily\Large\bfseries\color{ink} The rules are most of the studio}\\[8pt]
{\color{ink60} 183 of them, over 28 families, each able to stop a request and say
why. Not validation: the accumulated knowledge of what makes a map fail, written as
things the software will not let happen.}\\[9pt]

\begin{tabularx}{\linewidth}{L{52mm} Y}
\toprule
\textbf{Family} & \textbf{About} \\
\midrule
\id{SK} 28 & the sketch: shapes, layers, what drew no ground \\
\id{DR} 17 & dressing: where a thing may rest, what it may not close off \\
\id{PL} 16 & the plan: pieces, walls, spawns, objectives \\
\id{WX} 13 & the room stamped for a spawn or a wool \\
\id{OB} 12 & objectives: where a goal may stand, how a match can end \\
\id{CT} 10 & the middle: the crossing, its holes, its build zones \\
\bottomrule
\end{tabularx}

\vspace{12pt}
{\color{ink60}\small A rule is not a constant. Many read \emph{bands} whose values
depend on what the board is played for, so three boards of experience on one
objective teach nothing about the next.}
\end{minipage}
\end{minipage}

\vspace{10pt}
{\color{rule}\rule{\textwidth}{1pt}}\par\vspace{10pt}

% ------------------------------------------------------------- the two maps ---
{\sffamily\Large\bfseries\color{ink} Two boards made this way}\par\vspace{8pt}

% A board card: the render at the column's full width, its name and what it is.
% The two columns are narrow because an isometric render is nearly as tall as it is
% wide, and the sheet has one band left under the register.
\newcommand{\boardcard}[4]{%
  \begin{minipage}[t]{0.165\textwidth}
    \IfFileExists{../figures/generated/maps/#1-iso.png}{%
      \includegraphics[width=\linewidth]{../figures/generated/maps/#1-iso.png}}{%
      \fcolorbox{rule}{wash}{\parbox[c][150pt][c]{\dimexpr\linewidth-2\fboxsep-2\fboxrule\relax}%
        {\centering\color{ink30}\footnotesize board render}}}\par\vspace{7pt}
    {\sffamily\bfseries\color{ink} #2}\par\vspace{1pt}
    {\color{ink60} #3}\par\vspace{4pt}
    {\footnotesize\color{ink60} #4}
  \end{minipage}%
}

\noindent\begin{minipage}[t]{\textwidth}
\noindent
\boardcard{whitegape}{Whitegape}{destroy the monument, 20 a side}{Authored from
  nothing, to no brief. A limestone gorge quarried from both rims, each team's
  monument on the floor of its own pit.}
\hfill
\boardcard{hushwater}{Hushwater}{capture the wool, 18 a side}{Pulled off the
  composer at seed 5 and then adapted. The arrangement is the studio's; the place
  is the author's.}
\hfill
\begin{minipage}[t]{0.58\textwidth}
  {\sffamily\bfseries\color{ink} What a session looks like}\par\vspace{8pt}
  {\color{ink60}\small Every call each board's author made was recorded. An
  authoring run is not a single request: it is hundreds of reads against a handful
  of writes, because an agent cannot see the map and every belief it holds is a
  number the studio handed back.}\par\vspace{10pt}
  {\color{ink60}\small A 200 is not a yes. A request can succeed and still drop
  what it was sent, and the only place that is said is the \id{warnings} key. That
  is the most expensive lesson in the whole corpus.}\par\vspace{10pt}
  {\color{ink60}\small One recorded destroy-board run: 433 calls, 298 reads
  against 135 writes, 33 distinct routes, and 16 rule ids over 162 firings. Six of
  those were refusals. The other 156 arrived on a request that had already
  succeeded.}
\end{minipage}
\end{minipage}

\vspace{12pt}
{\color{rule}\rule{\textwidth}{1pt}}\par\vspace{6pt}
{\footnotesize\sffamily\color{ink60}
Every figure on this sheet was drawn by the studio being described.
The full paper, with both worked maps and the complete rule catalogue, is
\textit{How a PGM Map Gets Made}.}

\end{document}
